\section*{Introduction}

Accurate short-term knowledge of traffic conditions is essential for understanding and managing the dynamics of urban road networks.
However, traffic evolution is difficult to predict because it results from the interaction of two strongly coupled components.
On the one hand, traffic at a given location follows temporal patterns determined by its recent history; on the other hand,
changes occurring on one road segment can propagate through the network and affect other connected locations.
A forecasting model should therefore be able to exploit both temporal dependencies and the spatial structure of the road network.

An additional challenge arises from the fact that prediction quality is not uniform across operating conditions.
Models trained predominantly on regular traffic patterns may behave differently when confronted with less common situations,
such as widespread congestion.
In this context, evaluating only the average prediction error provides an incomplete picture of model performance.
A useful forecasting system should also provide an indication of its uncertainty and, importantly, this uncertainty
should remain informative when the input distribution differs from the conditions most frequently observed during training.


This project investigates graph-based probabilistic traffic forecasting under distribution shift using the METR-LA dataset,
which contains traffic-speed measurements collected from 207 sensors in the Los Angeles road network.
The forecasting problem is formulated using leakage-free temporal windows, with 12 historical observations used to predict
the following 12 time steps.
Performance is evaluated at three representative forecasting horizons: 15, 30, and 60 minutes ahead.
A chronological train-validation-test split is adopted throughout the experimental pipeline to preserve the temporal
structure of the data and prevent information leakage.


The study follows a progressive modeling approach.
A simple persistence baseline is first established as a reference for point forecasting performance.
A GRU-based temporal model without graph message passing is then introduced to capture temporal dependencies independently
for each sensor.
The main architecture extends this formulation with diffusion-based graph message passing, allowing information to
propagate across the road network before being processed by the recurrent temporal component.
Rather than producing only deterministic forecasts, both learned models estimate the 0.1, 0.5, and 0.9 conditional
quantiles through masked pinball loss, enabling the construction of predictive intervals and the joint evaluation of
accuracy and uncertainty.


Particular attention is devoted to understanding the contribution of the graph rather than considering it an inherently
beneficial architectural component.
The graph-temporal model is therefore compared with its no-graph counterpart and analyzed through architectural variants
and graph ablations, including different diffusion depths and weighted versus binarized adjacency.
The analysis further examines individual sensors and forecasting windows to identify conditions in which spatial message
passing improves or degrades predictions, and relates these effects to properties of the underlying graph topology.
Computational cost and model complexity are also considered alongside predictive performance and reliability.


Finally, the robustness of the forecasting system is evaluated under a predefined distribution shift based on high-congestion periods.
Congestion is identified using thresholds derived exclusively from the training data, allowing normal and stressed test
conditions to be compared without introducing test-set information into the shift definition.
Changes in point errors, prediction-interval coverage, interval width, and calibration are analyzed across forecasting horizons.
A calibration-oriented mitigation strategy is subsequently evaluated, together with an additional proposed approach for
improving robustness under congestion.
Overall, the project aims not only to determine how accurately the models forecast traffic, but also to investigate when
spatial information is useful, how uncertainty behaves when traffic conditions shift, and whether the resulting confidence
estimates remain reliable under challenging operating conditions.


